\documentclass[10pt,journal,compsoc]{IEEEtran}

\usepackage[utf8]{inputenc}
\usepackage[T1]{fontenc}
\usepackage{textcomp}
\ifCLASSOPTIONcompsoc
  \usepackage[nocompress]{cite}
\else
  \usepackage{cite}
\fi
\usepackage{amsmath,amssymb}
\usepackage{booktabs}
\usepackage{array}
\usepackage{graphicx}
\usepackage{xcolor}
\usepackage{pifont}
\usepackage{url}
\usepackage[hidelinks]{hyperref}
\def\UrlBreaks{\do\/\do-\do\_\do\.\do\:}
\Urlmuskip=0mu plus 1mu\relax
\renewcommand{\arraystretch}{1.1}

\title{GOV-AR: Risk-Bounded Multi-Tenant LLM Admission and Routing\\
under Delayed Token-Cost Feedback}

\author{Anonymous Authors}

\begin{document}

\IEEEtitleabstractindextext{%
\begin{abstract}
Shared enterprise LLM gateways face a control problem that post-hoc reporting
does not solve: output-token cost is unknown at admission time, several
requests may already be in flight, and provider telemetry may settle late. We
study this problem through GOV-AR, a governance-aware admission, routing, and
reservation path for multi-tenant AI gateways under delayed token-cost
feedback. The current artifact combines hard governance filtering,
tenant-isolated settled-versus-reserved budget accounting, request-level
reservation and idempotent settlement, and first-class \texttt{admit},
\texttt{queue}, \texttt{reject}, and \texttt{abstain} outcomes. The
implementation spans a reproducible research workspace, an operator-side
\texttt{gov-ar-admission} service, an optional PostgreSQL-backed ledger, and a
proxy-triggered live path. We report trace-driven evidence across delayed
settlement, multi-tenant contention, drift, fault handling, scalability, live
Azure validation, and ablation. The current results establish a reproducible
systems-and-method artifact, but not yet a full prompt-complete or
submission-complete Q1 study; we document the remaining gaps explicitly.
\end{abstract}

\begin{IEEEkeywords}
LLM gateways, admission control, delayed feedback, Kubernetes operators,
multi-tenant systems, FinOps, sovereign AI, Azure OpenAI.
\end{IEEEkeywords}}

\maketitle
\IEEEdisplaynontitleabstractindextext
\IEEEpeerreviewmaketitle

\section{Introduction}
Cost control for enterprise LLM gateways is difficult because the financial
commitment of a request is partly hidden until the response completes. A budget
controller that only tracks already settled usage ignores costs already engaged
by in-flight requests. At the opposite extreme, reserving for worst-case output
length is safe but wasteful. GOV-AR targets this middle ground: it decides
whether to admit, queue, reject, or abstain, selects a deployment, and commits
a reservation while respecting hard governance constraints.

The operator used in this work already contains telemetry ingestion, cost
attribution, budget tracking, sovereignty enforcement, route recommendation,
and Kubernetes-native policy surfaces. The missing piece is a reserve-dispatch-
settle control path for uncertain and delayed-cost requests. This article uses
the operator as the implementation substrate for that scientific decision
problem rather than as the article's end goal.

\section{Motivation and Failure Example}
Consider two tenants sharing the same premium deployment. Tenant A has already
spent most of its budget window. Tenant B has a fresh budget, but several long
responses are still in flight. If the controller admits a new request based
only on settled spend, it may over-commit. If it reserves at maximum output
length for every request, it may reject requests that would have remained safe
under calibrated reservations. This motivates a controller with explicit
liability tracking and conservative fallback when uncertainty grows.

\section{Related Work}
Adaptive LLM routing under budget constraints has been studied in recent work
such as Adaptive LLM Routing under Budget Constraints~\cite{panda2025adaptive}.
Cost-aware cascades and routers such as FrugalGPT~\cite{chen2023frugalgpt},
RouteLLM~\cite{ong2024routellm}, HYBRID LLM~\cite{ding2024hybrid}, and
RouteLMT~\cite{luo2026routelmt} show that query-dependent selection can improve
quality-cost trade-offs. Conformal routing adds distribution-free safety at the
routing level~\cite{uddin2026conformal}. Recent ACL 2026 work extends this
space toward broad model pools, length-aware routing, and multi-turn
trajectories~\cite{shi2026inferencedynamics,fu2026flare,zhang2026mtrouter}.
From the systems side, output-length-aware serving and unknown-length-aware
scheduling motivate reservation and length prediction
mechanisms~\cite{qiu2024ssjf,xie2026forelen,ruan2026libra}. Foundational theory
for delayed feedback and resource-constrained online decisions appears in
online learning under delayed feedback~\cite{joulani2013delayed} and contextual
bandits with knapsack-style constraints~\cite{badanidiyuru2014resourceful,han2023cbwk}.

The remaining gap is their combination with hard governance filtering,
tenant-isolated budget liability, explicit queue and abstain outcomes, and a
Kubernetes-native implementation path. The literature review is materially
improved, but still not saturated to the final target size and stopping rule
required by the original prompt.

\section{Problem Formulation}
For each request $r$ from tenant $t$, GOV-AR chooses one of
\texttt{admit(model,reservation)}, \texttt{queue}, \texttt{reject}, or
\texttt{abstain}. Before scoring, the candidate set is hard-filtered by allowed
provider, allowed zone, sensitivity compatibility, routability, approval state,
and model availability. The cost model is
\[
C_{t,m}=p_{in}(m)\cdot n_{in}(t)+p_{out}(m)\cdot N_{out}(t,m),
\]
where $N_{out}(t,m)$ is unknown at admission time.

The core accounting abstraction is
\[
\mathrm{available}_t(k)=\mathrm{budget}_t-\mathrm{settled}_t(k)-\mathrm{reserved}_t(k).
\]
This separates already observed spend from in-flight financial liability.

\section{GOV-AR Design}
The current design includes four coupled components.

\paragraph{Governance filter.}
Non-compliant models are eliminated before any soft ranking step.

\paragraph{Reservation policy.}
The present artifact supports empirical quantile reservation and mean-plus-
standard-deviation reservation over output-length estimates.

\paragraph{Ledger and settlement.}
The request ledger stores reservation amount, tenant, selected deployment,
pricing version, reservation mode, settlement status, and cancellation state.
Settlement is idempotent for repeated delivery of the same event identifier.

\paragraph{Conservative fallback.}
When drift or insufficient evidence invalidates a confident decision, the
controller can queue or abstain instead of silently continuing with stale
assumptions.

\section{Safety and Risk Analysis}
The strongest safety property currently evidenced is deterministic accounting
discipline for the implemented path: reservations are unique per request,
duplicate settlement identifiers are tolerated idempotently, conflicting
duplicate settlements are rejected, and cancellation releases reserved
liability. The current artifact does not yet provide a full prompt-grade
probabilistic guarantee with calibration intervals, coverage analysis by tenant,
and a final risk-bounded theorem. Scientific claims should therefore be read as
observed behavior plus partial design-level safety, not as a complete proven
guarantee.

\section{Kubernetes and Gateway Implementation}
The implementation is split between the research workspace in
\texttt{article3/} and operator-side code under \texttt{operateur/}. The
research workspace provides the admission kernel, ledger, predictors, replay
harness, experiments, provenance, and manuscript assets. The operator-side
slice adds:
\begin{itemize}
  \item \texttt{internal/govar/} for CRD-to-snapshot translation and decision
        backends;
  \item \texttt{cmd/gov-ar-admission} for \texttt{/v1/admit},
        \texttt{/v1/settle}, \texttt{/v1/cancel}, and
        \texttt{/v1/liability/\{tenant\}};
  \item an optional PostgreSQL-backed transactional ledger;
  \item chart templates to deploy the admission service;
  \item a proxy-triggered path in \texttt{header-proxy};
  \item webhook-sidecar propagation of GOV-AR configuration through pod
        annotations.
\end{itemize}
This is a meaningful operator integration slice, but it is still not a final
Envoy-native \texttt{ext\_proc} or external-authorization production path.

\section{Experimental Methodology}
The artifact contains a frozen protocol file, raw outputs, processed outputs,
and reproducible orchestration scripts. However, the current execution depth
must be distinguished from the ideal prompt protocol. The evidence package is
reproducible, but most experiment families remain below the requested scale in
seed count, request count, scenario breadth, and statistical completion. The
current repository now also includes an automatically generated statistical
summary over the available E1 and E2 campaign outputs, using paired bootstrap
intervals over run-level metric differences.

\section{Budget Safety and Utilization}
In the deterministic E1 summary, quantile reservation admitted two requests and
settled a total cost of 12.0 with 2.0 slack and zero overshoot. The
mean-plus-standard-deviation variant admitted one request, queued one, settled
7.0, and also had zero overshoot. In the campaign comparison, the trade-off
shifts: one variant can reduce overshoot while increasing slack. This supports
the central practical tension between utilization and conservatism. In the
current paired campaign analysis, the E1 quantile variant shows higher
overshoot than mean-plus-standard-deviation, with a mean overshoot difference
of 1.18 and a paired bootstrap 95\% interval of approximately
$[0.19, 2.35]$, while also using less slack by about 1.6 units.

\section{Multi-Tenant Routing}
In the deterministic E2 trace, both compared variants admitted three requests
and queued one, but the quantile variant incurred 0.5 overshoot while the
mean-plus-standard-deviation variant used more slack and avoided overshoot.
Tenant-level settled totals were recorded separately, confirming that the
current harness tracks tenant-isolated accounting. That said, the prompt-grade
multi-tenant study with six rich tenant profiles and large-scale scenarios is
still incomplete. In the available E2 campaign analysis, the sign of the trade
off reverses: quantile reservation reduces overshoot relative to
mean-plus-standard-deviation by about 1.62 units on average, with a paired
bootstrap 95\% interval of approximately $[-2.60, -0.59]$.

\section{Drift and Fault Injection}
The drift experiment demonstrates an observable conservative mode. After a
measured ratio drift above threshold, the artifact falls back to a safer
candidate and abstains when no sufficiently supported candidate remains. Fault
injection validates duplicate-settlement handling, expiry release, and
telemetry-insufficiency abstention. What is still missing is the full
20-scenario fault matrix with the repetition depth requested in the prompt.

\section{Scalability and Live Azure Validation}
The local scalability sweep replays synthetic traces up to 1000 requests and
measures replay throughput in the lightweight harness. These numbers should be
read only as local replay-engine observations. The Azure live validation
successfully discovered and invoked three real deployments spanning Azure
OpenAI and Foundry-served Mistral, with HTTP 200 responses and traceable token
usage or latency metadata. This is a real live validation, but still not the
full prompt-scale experiment of at least 300 successful requests per model
across multiple independent windows.

\section{Ablation and Sensitivity}
The ablation package compares reservation and queue-policy variants and shows
that aggressive under-reservation increases overshoot while disabling queueing
turns deferred decisions into direct refusals. This is useful mechanistic
evidence, but it still falls short of the full A0--A9 by scenario matrix
requested in the prompt.

\section{Threats to Validity}
The main threats are scale, dataset realism, and maturity bias. Synthetic
traces do not fully capture production prompt distributions. Replay throughput
does not equal end-to-end gateway throughput. The literature review is stronger
than the initial state but not yet fully saturated. The manuscript has moved
closer to a scientific structure, yet it is still not a final Q1 submission.

\section{Discussion}
The strongest contribution of the current repository state is not a polished
paper claim, but a hardening of the implementation substrate: operator audit,
reserve-settle scaffold, operator-side admission service, optional PostgreSQL
ledger, local Kind automation, live Azure validation, and a now much clearer
prompt-compliance audit trail. The main remaining work is therefore not finding
small wording tweaks, but closing the genuine scientific and experimental gaps.

\section{Conclusion}
GOV-AR frames enterprise LLM routing as a joint admission, reservation, and
settlement problem under delayed token-cost feedback and hard governance
constraints. The present artifact already demonstrates a reproducible decision
kernel, a Kubernetes-oriented implementation path, local Docker/Kind/Helm
execution, and real Azure validation against existing deployments. It does not
yet satisfy the full original prompt, and it is not yet a finished Q1 article.
The next mandatory steps are full prompt-scale experimentation, stronger
statistical analysis, deeper Envoy-native integration, complete evidence
tracking, and final publication packaging.

\bibliographystyle{IEEEtran}
\bibliography{references}
\end{document}
